\documentclass[11pt,a4paper]{article}
\usepackage[T1]{fontenc}
\usepackage{amsmath,amssymb,amsthm}
\usepackage[margin=25mm]{geometry}
\usepackage{longtable,booktabs,array,calc}
\usepackage{needspace}
\usepackage{xurl}
\usepackage[hidelinks]{hyperref}
\hypersetup{pdftitle={Parity Statistics of Nested Floor
Powers: Conditional Descent Results for the Juggler
Map},pdfauthor={Philippe
Cochin},pdfsubject={Conditional mathematics preprint}}
\setlength{\emergencystretch}{2em}
\setlength{\tabcolsep}{4pt}
\renewcommand{\arraystretch}{1.12}
\setlength{\LTpre}{0.8em}
\setlength{\LTpost}{0.8em}
\clubpenalty=10000
\widowpenalty=10000
\allowdisplaybreaks[1]
\raggedbottom
\providecommand{\tightlist}{\setlength{\itemsep}{0pt}\setlength{\parskip}{0pt}}
\providecommand{\passthrough}[1]{#1}
\setcounter{secnumdepth}{-1}
\title{Parity Statistics of Nested Floor
Powers\\[0.4em]{\large Conditional Descent Results for the Juggler Map}}
\author{Philippe Cochin\\\small\texttt{philippe@cochin.fr}}
\date{9 September 2026}
\begin{document}
\maketitle
\begin{center}\small Preprint\end{center}
\section{Abstract}\label{abstract}

The Juggler map applies the integer part of the square root at even
positive integers and of the three-halves power at odd positive
integers. Its finite itineraries lead to joint parity questions for
compositions of floor powers. We give a self-contained conditional
framework relating these parity statistics to descent below the initial
value. Exact branch-indicator identities and a power envelope reduce
finite descent certificates to explicit correlation estimates along
prescribed formal chains. The classical second-derivative estimate gives
an unconditional certificate class of natural density \(3/4\). Under
stated finite-chain correlation hypotheses, the certificate classes of
lengths at most four and five have densities \(13/16\) and \(7/8\),
respectively. If every fixed-depth odd-rooted itinerary has its
fair-share density, the set of starts admitting some finite
power-envelope certificate has density one; no uniformity in depth is
needed for this implication. We also record exact floor-defect
identities, a Fourier criterion sufficient for the correlation
hypotheses, and a shift-averaged cancellation estimate. The nested
correlation and deterministic kernel estimates remain unproved here. In
particular, the larger certificate densities and the density-one
conclusion are conditional, and none asserts universal arrival at \(1\).

\textbf{Keywords:} Juggler map; nested floor powers; parity
correlations; discrepancy; stopping time; conditional descent.

\section{1. Scope and relation to earlier
work}\label{scope-and-relation-to-earlier-work}

For a positive integer \(n\), define \[
J(n)=
\begin{cases}
\lfloor n^{1/2}\rfloor,&n\ \text{even},\\
\lfloor n^{3/2}\rfloor,&n\ \text{odd}.
\end{cases}
\] This is the Juggler map recorded in OEIS A094683 {[}1{]}. The
conjecture that every positive orbit reaches \(1\) remains open. The
question studied here is more limited: which distribution estimates for
finite itineraries would imply that many starting values eventually fall
below their initial value?

Single-floor parity estimates follow from classical exponential-sum
methods {[}2, 3{]}. Composing floors creates an additional difficulty.
For \(m=\lfloor n^{3/2}\rfloor\) and \(\theta=\{n^{3/2}\}\), \[
m^{3/2}=n^{9/4}-\tfrac32\theta n^{3/4}+O(n^{-3/4}).
\] The coefficient of the fractional part grows with \(n\). A theorem
for the smooth phase \(n^{9/4}\) therefore does not by itself estimate
the nested phase. Nor does a parity estimate over consecutive starting
values automatically apply to a sparse set of orbit images.

The contribution of this note is an explicit reduction: the exact
identities, finite certificate classification, sufficient analytic
hypotheses, and conditional consequences are stated separately. The
single-floor estimate and the probabilistic counting argument use
standard methods. The latter is in the spirit of stopping-time arguments
for the Collatz map, such as Terras {[}4{]}; no theorem about Collatz is
transferred to the Juggler map.

This version replaces an earlier working draft that asserted nested
parity equidistribution and kernel cancellation. Those analytic claims
are not established by the arguments retained here. The powers \(23/24\)
and \(95/96\) below are research targets, not proved error exponents.
The former numerical threshold \(3.6\cdot10^{13}\) is not a certified
threshold for any nested cancellation theorem in this version. All
proofs of conditional implications are given below; their analytic
hypotheses remain explicit.

Throughout, \(e(t)=\exp(2\pi it)\), \(\{t\}=t-\lfloor t\rfloor\), and
\(\psi(t)=(-1)^{\lfloor t\rfloor}\). Natural density is measured among
all positive integers. Thus a length-\(d\) odd-rooted word with count
\(2^{-d}N+o(N)\) has relative density \(2^{-(d-1)}\) among odd starts.
Constants in \(O\) and \(\ll\) may depend on fixed words and explicit
fixed parameters. They are not asserted to be uniform in depth.

\section{2. Exact itineraries and the power
envelope}\label{exact-itineraries-and-the-power-envelope}

Let \(w=w_1\cdots w_d\) be a word in \(\{E,O\}\). The itinerary
\(\operatorname{word}_d(n)\) records the parities of
\(n,J(n),\ldots,J^{d-1}(n)\). It contains \(d\) letters and specifies
the \(d\) operations that produce \(J^d(n)\). Write \(o(w)\) for the
number of its odd letters.

\needspace{5\baselineskip}

\textbf{Proposition 2.1 (power envelope).} If \(w\) is realized at
\(n\), then \[
\bigl(J^d(n)\bigr)^{2^d}\le n^{3^{o(w)}}.
\] Consequently, if \(n\ge2\) and \(3^{o(w)}<2^d\), then \(J^d(n)<n\).

\emph{Proof.} At any positive integer \(x\), the even operation
satisfies \(J(x)^2\le x\), and the odd operation satisfies
\(J(x)^2\le x^3\). Suppose a prefix of length \(r\) with \(o\) odd
letters ends at \(x\) and obeys \(x^{2^r}\le n^{3^o}\). An even
operation preserves the right-hand exponent after raising
\(J(x)^2\le x\) to \(2^r\); an odd operation replaces it by \(3^{o+1}\).
Induction starts at the empty prefix. If \(3^{o(w)}<2^d\) and \(n>1\),
the bound is strictly less than \(n^{2^d}\), proving descent.
\(\square\)

A word satisfying \(3^{o(w)}<2^{|w|}\) is called \emph{contracting}. A
realized contracting prefix is a \emph{power-envelope descent
certificate}. Let \[
\mathcal C_d=\{n\ge2:\text{some contracting prefix of length at most }d
\text{ is realized at }n\}.
\] Also set \[
\mathcal C_\infty=\bigcup_{d\ge1}\mathcal C_d.
\] These sets refer to this particular sufficient certificate. Flooring
can cause descent even when the envelope does not certify it. Therefore
\(\mathcal C_d\) need not equal the set of all starts that descend
within \(d\) steps.

For an odd-rooted target word \(w\), define its \emph{formal chain} for
every odd \(n\), regardless of whether \(n\) realizes \(w\): \[
x_1^w(n)=n,\qquad
\xi_t^w(n)=\bigl(x_t^w(n)\bigr)^{p_t},\qquad
x_{t+1}^w(n)=\lfloor\xi_t^w(n)\rfloor\quad(1\le t<d),
\] where \(p_t=3/2\) for \(w_t=O\) and \(p_t=1/2\) for \(w_t=E\). Set
\(s_t^w(n)=\psi(\xi_t^w(n))\), and put \(\epsilon_t=1\) for
\(w_{t+1}=E\), \(\epsilon_t=-1\) for \(w_{t+1}=O\).

\needspace{5\baselineskip}

\textbf{Lemma 2.2 (parity and branch indicators).} For real \(x\),
\(\lfloor x\rfloor\) is odd exactly when \(\{x/2\}\in[1/2,1)\). For
every odd \(n\) and odd-rooted word \(w\), \[
\mathbf1_{\operatorname{word}_d(n)=w}
=2^{-(d-1)}\prod_{t=1}^{d-1}(1+\epsilon_t s_t^w(n)).
\tag{2.1}
\]

\emph{Proof.} Writing \(x=\lfloor x\rfloor+\{x\}\) proves the first
assertion, including the half-open endpoints. Each factor divided by two
tests the parity of \(x_{t+1}^w\). Until the first failed test, the
formal chain agrees with the true orbit, by induction on \(t\). A failed
test makes the product zero and prevents realization of \(w\). If all
tests pass, the chains agree through the required letters and the
product is one. The empty product covers \(d=1\). \(\square\)

The use of formal chains is essential: the product is defined on every
odd input before any correlation or Fourier expansion is made.

\section{3. An unconditional certificate
density}\label{an-unconditional-certificate-density}

Let \(M(N)=\#\{n\le N:n\text{ odd}\}=N/2+O(1)\), and set \[
S_O(N)=\sum_{\substack{n\le N\\n\ \mathrm{odd}}}\psi(n^{3/2}).
\]

\needspace{5\baselineskip}

\textbf{Theorem 3.1 (single-floor parity).} One has \[
S_O(N)=O(N^{5/6}).
\] In particular, \[
\#\{n\le N:\operatorname{word}_2(n)=OO\}
=N/4+O(N^{5/6}),
\] \[
\#(\mathcal C_2\cap[1,N])=3N/4+O(N^{5/6}).
\]

\emph{Proof.} Write \(n=2r+1\) and \(g(r)=\tfrac12(2r+1)^{3/2}\). Then
\(g''(r)=\tfrac32(2r+1)^{-1/2}\). On a dyadic block with \(r\asymp Q\),
the classical second-derivative test {[}2{]} gives \[
\left|\sum_{r\asymp Q}e(hg(r))\right|
\ll h^{1/2}Q^{3/4}+h^{-1/2}Q^{1/4}
\quad(h\ge1).
\] The Erd\H{o}s--Tur\'{a}n inequality {[}3, Chapter 2{]} bounds the
unnormalized interval discrepancy by \[
\ll Q/H+\sum_{h=1}^{H}\frac1h
\left|\sum_{r\asymp Q}e(hg(r))\right|
\ll Q/H+H^{1/2}Q^{3/4}+Q^{1/4}.
\] Taking \(H=\lfloor Q^{1/6}\rfloor\) for sufficiently large \(Q\)
gives \(O(Q^{5/6})\). Lemma 2.2 identifies parity with membership in
\([1/2,1)\); summing the dyadic bounds proves the first assertion. Since
\(S_O=M-2\#OO\), it also proves the count of \(OO\).

The only minimal contracting words of length at most two are \(E\) and
\(OE\). The first class has count \(\lfloor N/2\rfloor\), and the second
has count \(M-\#OO=N/4+O(N^{5/6})\). They are disjoint, and Proposition
2.1 proves the stated certificate count. \(\square\)

This argument estimates starting values in an interval. It makes no
independence assertion about successive orbit parities.

\section{4. Explicit correlation
hypotheses}\label{explicit-correlation-hypotheses}

For a fixed odd-rooted word \(w\) of length \(d\) and nonempty
\(A\subseteq\{1,\ldots,d-1\}\), define \[
R_{w,A}(N)=\sum_{\substack{n\le N\\n\ \mathrm{odd}}}
\prod_{t\in A}s_t^w(n).
\]

\needspace{5\baselineskip}

\textbf{Hypothesis \(\mathrm H(w;\delta)\).} For some fixed
\(0<\delta<1\), every nonempty \(A\) satisfies
\(R_{w,A}(N)=O_{w,A}(N^{1-\delta})\) as \(N\to\infty\).

\needspace{5\baselineskip}

\textbf{Hypothesis \(\mathrm H_0(w)\).} The same correlations satisfy
\(R_{w,A}(N)=o(N)\), without a prescribed rate.

These are assertions about the formal chains, not assumptions that the
true orbit is a sequence of independent coin tosses. Hypothesis
\(\mathrm H(w;\delta)\) implies \(\mathrm H_0(w)\).

\needspace{5\baselineskip}

\textbf{Proposition 4.1 (conditional word count).} Under
\(\mathrm H(w;\delta)\), \[
\#\{n\le N:\operatorname{word}_d(n)=w\}
=2^{-d}N+O_w(N^{1-\delta}).
\tag{4.1}
\] Under \(\mathrm H_0(w)\), the error is \(o(N)\).

\emph{Proof.} Expand (2.1) and sum. The empty subset contributes
\(2^{-(d-1)}M(N)=2^{-d}N+O(1)\). The other \(2^{d-1}-1\) terms are fixed
signed multiples of \(R_{w,A}\), and there are finitely many for fixed
\(w\). \(\square\)

For example, \(\mathrm H(OOEE;\delta)\) concerns the seven nonempty
products of \[
\psi(n^{3/2}),\qquad \psi(m^{3/2}),\qquad \psi(v^{1/2}),
\quad m=\lfloor n^{3/2}\rfloor,\quad v=\lfloor m^{3/2}\rfloor.
\] Proposition 4.1 then gives the count \(N/16+O(N^{1-\delta})\) needed
for four-step certificates. Theorem 3.1 establishes only the first of
these seven estimates; the others are not supplied here.

\subsection{4.1 A sufficient Fourier
criterion}\label{a-sufficient-fourier-criterion}

For \(\mathbf k=(k_1,\ldots,k_{d-1})\in\mathbb Z^{d-1}\), let \[
S_{w,\mathbf k}(N)=\sum_{\substack{n\le N\\n\ \mathrm{odd}}}
e\left(\frac12\sum_{t=1}^{d-1}k_t\xi_t^w(n)\right).
\]

\needspace{5\baselineskip}

\textbf{Proposition 4.2 (conditional Fourier-to-parity transfer).} Fix
\(d\ge2\), \(w\), and \(0<\eta,\sigma<1\). Suppose, uniformly over all
nonzero integer vectors with
\(\|\mathbf k\|_\infty\le\lfloor N^\eta\rfloor\), that \[
|S_{w,\mathbf k}(N)|\ll_w N^{1-\sigma}.
\tag{4.2}
\] Then \(\mathrm H(w;\delta)\) holds for every
\(0<\delta<\min(\eta,\sigma)\). The constants may depend on these fixed
parameters.

\emph{Proof.} Apply the multidimensional Erd\H{o}s--Tur\'{a}n--Koksma
inequality {[}3, Chapter 2, p.~116{]} to the points
\(\mathbf y_n=(\{\xi_1^w(n)/2\},\ldots,\{\xi_{d-1}^w(n)/2\})\). For any
half-open axis-parallel box \(B\subseteq[0,1)^{d-1}\), its unnormalized
discrepancy is \[
\begin{split}
\left|\#\{n\le N\text{ odd}:\mathbf y_n\in B\}-M(N)|B|\right|
&\ll_d \frac{M(N)}{H}
+\sum_{0<\|\mathbf k\|_\infty\le H}
\frac{|S_{w,\mathbf k}(N)|}{r(\mathbf k)},\\
r(\mathbf k)&=\prod_{t=1}^{d-1}\max(1,|k_t|).
\end{split}
\] The weight sum is \(O_d((1+\log H)^{d-1})\). With
\(H=\lfloor N^\eta\rfloor\), (4.2) gives box discrepancy
\(O_w(N^{1-\eta}+N^{1-\sigma}(1+\log N)^{d-1})\). Each nonempty sign
product is constant, with value \(1\) or \(-1\), on the \(2^{d-1}\)
boxes obtained by splitting each coordinate at \(1/2\). Its integral
over the unit cube is zero. Summing their discrepancies therefore bounds
\(R_{w,A}\); absorbing the fixed logarithmic power proves the assertion.
\(\square\)

The criterion needs every indicated mixed mode, including vectors with
zero coordinates and mixed signs. A bound for one undecorated kernel is
not a substitute for (4.2). Qualitatively, cancellation
\(S_{w,\mathbf k}(N)=o(N)\) for each fixed nonzero \(\mathbf k\) also
implies \(\mathrm H_0(w)\): keep \(H\) fixed in the displayed
inequality, let \(N\to\infty\), and then let \(H\to\infty\).

\section{5. Conditional finite-depth descent
densities}\label{conditional-finite-depth-descent-densities}

\needspace{5\baselineskip}

\textbf{Lemma 5.1 (minimal certificates through length five).} The
contracting words of length at most five with no proper contracting
prefix are exactly \[
E,\qquad OE,\qquad OOEE,\qquad OOOEE,\qquad OOEOE.
\tag{5.1}
\]

\emph{Proof.} An \(E\)-rooted word already contracts at its first
letter, and an \(OE\)-rooted word contracts at its second. Every
remaining word begins with \(OO\). At length three, even two odd letters
give \(3^2>2^3\), so no new certificate appears. At length four,
contraction requires at most two odd letters, giving only \(OOEE\).
After excluding that prefix, the possible fourth-level prefixes are
\(OOEO\), \(OOOE\), and \(OOOO\). At length five, contraction requires
at most three odd letters because \(3^3<2^5<3^4\). The first two
prefixes must therefore end with \(E\), and the third cannot contract.
This proves (5.1). \(\square\)

\needspace{5\baselineskip}

\textbf{Theorem 5.2 (conditional four-step density).} If
\(\mathrm H(OOEE;\delta_4)\) holds for some \(\delta_4>0\), then \[
\#(\mathcal C_4\cap[1,N])
=\frac{13N}{16}+O(N^{1-\delta_*}),
\qquad \delta_* =\min(1/6,\delta_4).
\] The qualitative hypothesis \(\mathrm H_0(OOEE)\) instead suffices for
natural density \(13/16\).

\emph{Proof.} Lemma 5.1 gives the disjoint decomposition into the prefix
classes \(E\), \(OE\), and \(OOEE\). Theorem 3.1 supplies the first two
counts, and Proposition 4.1 supplies the third. Their densities add to
\(1/2+1/4+1/16=13/16\). \(\square\)

\needspace{5\baselineskip}

\textbf{Theorem 5.3 (conditional five-step density).} Suppose
\(\mathrm H(w;\delta_w)\) holds for each \[
w\in\mathcal W=\{OOEE,OOOEE,OOEOE\},\qquad \delta_w>0.
\] Then \[
\#(\mathcal C_5\cap[1,N])=\frac{7N}{8}+O(N^{1-\delta_*}),
\qquad \delta_* =\min\bigl(1/6,\min_{w\in\mathcal W}\delta_w\bigr).
\] The three qualitative hypotheses \(\mathrm H_0(w)\) suffice for
natural density \(7/8\).

\emph{Proof.} The two additional prefix classes in (5.1) are disjoint
from each other and from the four-step classes. Each has density
\(1/32\) by Proposition 4.1. Thus the total is \(13/16+1/32+1/32=7/8\).
Lemma 5.1 proves that this list exhausts power-envelope certificates
through length five. \(\square\)

The exact fractions and their analytic requirements can be read
together:

\begingroup\small

\begin{longtable}[]{@{}
  >{\raggedright\arraybackslash}p{(\linewidth - 6\tabcolsep) * \real{0.2500}}
  >{\raggedright\arraybackslash}p{(\linewidth - 6\tabcolsep) * \real{0.2500}}
  >{\raggedright\arraybackslash}p{(\linewidth - 6\tabcolsep) * \real{0.2500}}
  >{\raggedright\arraybackslash}p{(\linewidth - 6\tabcolsep) * \real{0.2500}}@{}}
\toprule\noalign{}
\begin{minipage}[b]{\linewidth}\raggedright
Certificate length
\end{minipage} & \begin{minipage}[b]{\linewidth}\raggedright
Minimal prefixes added
\end{minipage} & \begin{minipage}[b]{\linewidth}\raggedright
Cumulative density
\end{minipage} & \begin{minipage}[b]{\linewidth}\raggedright
Status
\end{minipage} \\
\midrule\noalign{}
\endhead
\bottomrule\noalign{}
\endlastfoot
1 & \(E\) & \(1/2\) & Unconditional \\
2 & \(OE\) & \(3/4\) & Theorem 3.1 \\
4 & \(OOEE\) & \(13/16\) & Conditional on \(\mathrm H_0(OOEE)\) \\
5 & \(OOOEE,OOEOE\) & \(7/8\) & Conditional on the three hypotheses in
Theorem 5.3 \\
\end{longtable}

\endgroup

For the five-letter words, the four tested real images are,
respectively, \((n^{3/2},m^{3/2},v^{3/2},z^{1/2})\), with
\(z=\lfloor v^{3/2}\rfloor\), and \((n^{3/2},m^{3/2},v^{1/2},u^{3/2})\),
with \(u=\lfloor v^{1/2}\rfloor\). Each hypothesis concerns all fifteen
nonempty sign products of its four coordinates.

Separately, if \(\mathrm H_0(w)\) holds for every odd-rooted word of
length four, all eight such classes have density \(1/16\). This complete
census is conditional as well; it is stronger than the single
four-letter hypothesis used in Theorem 5.2.

\section{6. The fixed-depth reduction to density-one
certificates}\label{the-fixed-depth-reduction-to-density-one-certificates}

\needspace{5\baselineskip}

\textbf{Hypothesis \(\mathrm{FD}\).} For every fixed \(d\ge1\) and every
odd-rooted word \(w\) of length \(d\), \[
\#\{n\le N:\operatorname{word}_d(n)=w\}=2^{-d}N+o_w(N).
\] No common error rate in \(d\) is assumed. The family
\(\mathrm H_0(w)\) over all such words implies \(\mathrm{FD}\).

\needspace{5\baselineskip}

\textbf{Theorem 6.1 (conditional density-one certificate theorem).}
Under \(\mathrm{FD}\), \(\mathcal C_\infty\) has natural density one.

\emph{Proof.} Put \(p=\log2/\log3>1/2\). A word with no contracting
prefix through length \(d\) must in particular have at least \(pd\) odd
letters at length \(d\). All such words begin with \(O\). For a
uniformly chosen odd-rooted word of length \(d\), the number of odd
letters is distributed as \(1+B_{d-1}\), where \(B_{d-1}\) is a binomial
random variable with parameters \((d-1,1/2)\). This is an auxiliary
count on words, not a stochastic model asserted for the Juggler orbit.

Choose \(q=(p+1/2)/2\), so \(1/2<q<p\). For all sufficiently large fixed
\(d\), \(pd-1\ge q(d-1)\). For \(t>0\), the exponential Markov
inequality and the binomial generating function give \[
\Pr(B_{d-1}\ge q(d-1))
\le\left(\frac{1+e^t}{2e^{qt}}\right)^{d-1}.
\] The logarithm of the expression in parentheses is zero at \(t=0\) and
has derivative \(1/2-q<0\) there. Hence some fixed \(t>0\) makes it a
number \(\vartheta<1\).

For fixed \(d\), \(\mathrm{FD}\) and a finite sum over surviving words
show that the natural density of \(\mathbb N\setminus\mathcal C_d\) is
their number divided by \(2^d\). It is at most
\(\tfrac12\vartheta^{d-1}\) for sufficiently large \(d\). Since
\(\mathcal C_d\subseteq\mathcal C_\infty\), \[
\overline{\operatorname{dens}}(\mathbb N\setminus\mathcal C_\infty)
\le\tfrac12\vartheta^{d-1}.
\] Letting \(d\to\infty\) proves the assertion. In this argument
\(N\to\infty\) is taken first, with \(d\) fixed. \(\square\)

The conclusion concerns descent below the initial value. It does not
prove that a descended value reaches \(1\), or that subsequent orbit
samples avoid an exceptional set. Even a density-zero exceptional set
can contain a nontrivial cycle. No interchange of the two limits, and no
assertion of finite stopping time for every start, follows from Theorem
6.1.

\section{7. Exact floor defects and the unresolved analytic
estimates}\label{exact-floor-defects-and-the-unresolved-analytic-estimates}

For \(n\ge3\), write \[
X=n^{3/2},\quad m=\lfloor X\rfloor,\quad\theta=X-m,
\qquad Y=m^{3/2},\quad v=\lfloor Y\rfloor,\quad\theta_2=Y-v.
\]

\needspace{5\baselineskip}

\textbf{Lemma 7.1 (one-signed linearization).} One has \[
m^{3/2}=\tfrac32mn^{3/4}-\tfrac12n^{9/4}+E(n),
\qquad
0\le E(n)\le\tfrac38(n^{3/2}-1)^{-1/2}\le\tfrac12n^{-3/4}.
\]

\emph{Proof.} Let \(a=\sqrt m\), \(b=\sqrt X\). Direct factorization
gives \[
E=a^3-\tfrac32a^2b+\tfrac12b^3
=\tfrac12(a-b)^2(2a+b)\ge0.
\] Moreover \(\theta=b^2-a^2\), and \(Ea\le\tfrac38\theta^2\) follows
from \(4a(2a+b)\le3(a+b)^2\), equivalent to \((5a+3b)(a-b)\le0\). Thus
\(E\le\tfrac38m^{-1/2}\), which is at most \(\tfrac38(X-1)^{-1/2}\).
Finally \((X-1)^{-1/2}\le\tfrac43X^{-1/2}\) for \(n\ge3\). \(\square\)

\needspace{5\baselineskip}

\textbf{Lemma 7.2 (gap and carry).} For an integer \(h\ge1\), set
\(\Delta_hX(n)=X(n+2h)-X(n)\). Then \[
m(n+2h)-m(n)=\lfloor\Delta_hX(n)\rfloor+\kappa_h(n),
\quad
\kappa_h(n)=\mathbf1_{\{X(n)\}+\{\Delta_hX(n)\}\ge1}\in\{0,1\}.
\]

\emph{Proof.} Expand \(X(n)+\Delta_hX(n)\) into integer and fractional
parts, and take its floor. The sum of the two fractional parts lies in
\([0,2)\). \(\square\)

For \(1\le h\le P/4\) and \(n\in(P,2P]\), differentiation gives
\((\Delta_hX)'(n)\asymp hP^{-1/2}\). The level sets of
\(\lfloor\Delta_hX\rfloor\) consequently form \(O(1+hP^{1/2})\)
intervals. Full intervals have length \(\asymp P^{1/2}/h\); the two
endpoint intervals can be shorter. A coefficient involving this floor is
frozen only on these intervals. A derivative estimate proved there must
be summed over the partition, or replaced by a separate global argument
with controlled errors.

\needspace{5\baselineskip}

\textbf{Lemma 7.3 (next-level defect).} One has \[
v^{3/2}=m^{9/4}-\tfrac32v^{1/2}\theta_2-\mathcal R,
\qquad 0\le\mathcal R\le\tfrac38v^{-1/2}\theta_2^2.
\]

\emph{Proof.} Apply Taylor's theorem to \(x^{3/2}\) between \(v\) and
\(v+\theta_2=Y\). Its second derivative is
\(\tfrac34x^{-1/2}\le\tfrac34v^{-1/2}\) on that interval, and is
positive. Rearranging the Taylor formula gives the result. \(\square\)

These identities describe the objects requiring cancellation; they do
not supply that cancellation. Two targets from the earlier draft are: \[
\left|\sum_{\substack{P<n\le2P\\n\ \mathrm{odd}}}
e\bigl(\tfrac i2 n^{3/2}+\tfrac j2m^{3/2}\bigr)\right|
\ll_\varepsilon P^{23/24+\varepsilon},
\quad |i|\le2P^{1/24},\quad1\le|j|\le2P^{1/24},
\tag{7.1}
\] and, with \(c_k(n)=\tfrac{3k}{4}n^{9/8}\), \[
K_k(P)=\sum_{\substack{P<n\le2P\\n\ \mathrm{odd}}}
e(c_k(n)\theta_2(n)),
\qquad |K_k(P)|\ll_\varepsilon P^{95/96+\varepsilon},
\quad1\le k\le P^{1/24}.
\tag{7.2}
\] The indices \(i,j,k\) are integers. Both estimates, with the printed
uniformities, are \textbf{open targets in this note}. Equations (7.1)
and (7.2) are not invoked in any unconditional proof above. Even if
(7.2) were established, its transfer to the full mixed-mode families in
Proposition 4.2 would still require proof. In particular, a coefficient
\(j/2\) can be a half-integer, so a lemma restricted to integer
frequencies cannot be used without extending its hypotheses.

Short-interval variants require separate estimates. A bound for a
transition set \(\Omega\subset(P,2P]\) of the form \(|\Omega|\ll P^a\)
implies only \(|\Omega\cap I|\le\min(|I|,O(P^a))\). It does not imply a
bound proportional to \(|I|/P\). Therefore the earlier localized-kernel
threshold \(|I|\ge P^{29/48+\delta}\) is not asserted here.

\subsection{7.1 An unconditional estimate averaged over a
shift}\label{an-unconditional-estimate-averaged-over-a-shift}

\needspace{5\baselineskip}

\textbf{Proposition 7.4 (shift average).} Let \(L\ge1\), let
\(A_1<\cdots<A_L\) obey \(|A_t-A_s|\ge a|t-s|\) with \(a>0\), and let
\(x_1,\ldots,x_L\) be real. For
\(S_\lambda=\sum_{t=1}^{L}e(A_t\{x_t+\lambda\})\), \[
\left|\int_0^1|S_\lambda|^2\,d\lambda-L\right|
\le\frac4\pi\frac La(1+\log L).
\tag{7.3}
\] For every \(0<\eta<1\), outside a set of shifts of measure at most
\(\eta\), \[
|S_\lambda|\le
\left[\frac L\eta\left(1+\frac4{\pi a}(1+\log L)\right)\right]^{1/2}.
\]

\emph{Proof.} Expanding the square gives the diagonal term \(L\). For
\(t\ne s\), the integrand \(e(A_t\{x_t+\lambda\}-A_s\{x_s+\lambda\})\)
is periodic in \(\lambda\). Integrate over a period starting at one of
its two fractional-part breakpoints. There is at most one other interior
breakpoint, giving at most two intervals on which the phase is affine
with slope \(A_t-A_s\). On each interval the integral has modulus at
most \(1/(\pi|A_t-A_s|)\). Hence the off-diagonal contribution is
bounded by \[
\frac{2}{\pi a}\sum_{t\ne s}\frac1{|t-s|}
=\frac4{\pi a}\sum_{r=1}^{L-1}\frac{L-r}{r}
\le\frac4\pi\frac La(1+\log L).
\] For \(L=1\) the off-diagonal sum is empty. Markov's inequality
applied to \(|S_\lambda|^2\) gives the second assertion. \(\square\)

The estimate allows a small exceptional set of shifts. It gives no bound
for a specified shift, including \(\lambda=0\) in (7.2).

\section{8. Evidence, availability, and remaining
work}\label{evidence-availability-and-remaining-work}

The status of the principal statements is as follows.

\begingroup\small

\begin{longtable}[]{@{}
  >{\raggedright\arraybackslash}p{(\linewidth - 2\tabcolsep) * \real{0.5000}}
  >{\raggedright\arraybackslash}p{(\linewidth - 2\tabcolsep) * \real{0.5000}}@{}}
\toprule\noalign{}
\begin{minipage}[b]{\linewidth}\raggedright
Statement
\end{minipage} & \begin{minipage}[b]{\linewidth}\raggedright
Warrant in this version
\end{minipage} \\
\midrule\noalign{}
\endhead
\bottomrule\noalign{}
\endlastfoot
Envelope, branch identity, minimal words, floor defects & Exact proofs
in Sections 2, 5, and 7 \\
Certificate density \(3/4\) & Classical analytic proof in Theorem 3.1 \\
Correlation-to-count and Fourier-to-parity implications & Proofs in
Propositions 4.1 and 4.2 \\
Certificate densities \(13/16\), \(7/8\) & Conditional theorems; nested
hypotheses unproved \\
Density-one finite certificates & Conditional on \(\mathrm{FD}\); no
depth uniformity assumed \\
Deterministic nested sums and short-interval kernels & Unproved; no
numerical threshold certified \\
Shift-averaged estimate & Direct integral proof in Proposition 7.4 \\
\end{longtable}

\endgroup

Supporting software is available in the Balanced Ternary Mathematical
Laboratory repository, \url{https://github.com/sneakyweasel/btlab}. The
accompanying source package contains this manuscript, its LaTeX build,
and an exact-arithmetic validation script. That script checks the finite
certificate list, the indicator identity on a finite census, and
selected algebraic inequalities. Those checks support reproducibility;
they are not proofs of asymptotic cancellation.

Earlier repository audits and Lean modules concern identities,
inequalities, and arithmetic from the previous draft. This version
claims no Lean formalization of its analytic hypotheses or of the
complete conditional manuscript. Earlier theorem numbers are superseded
and should not be used as citations to unconditional nested estimates.
Companion texts that rely on those estimates require their own review.

To obtain an unconditional improvement beyond the certificate density
\(3/4\), it would suffice to establish the correlations identified in
Theorem 5.2. To reach \(7/8\) by this route requires the additional
five-letter correlations in Theorem 5.3. A proof of all fixed-depth
counts would imply density-one certificates by Theorem 6.1, but a
further argument would still be required to settle universal
termination.

\needspace{12\baselineskip}

\section{Acknowledgments and AI
assistance}\label{acknowledgments-and-ai-assistance}

Large language models were used throughout the development of the
earlier draft and this revision, including the formulation and review of
proof arguments, drafting, programming, and discussion of Lean
formalizations. The September 2026 review with OpenAI Codex identified
unresolved analytic steps and prepared this conditional formulation, its
proofs, and its release materials. AI assistance and automated checks
are not independent mathematical validation. The author is responsible
for the statements, proofs, code, and final approval of the preprint.
The models are not authors.

\section{References}\label{references}

\begin{enumerate}
\def\labelenumi{\arabic{enumi}.}
\tightlist
\item
  OEIS Foundation Inc., ``Juggler sequence: if n mod 2 = 0 then
  floor(sqrt(n)) else floor(n\^{}(3/2)),'' \emph{The On-Line
  Encyclopedia of Integer Sequences}, A094683.
  \href{https://oeis.org/A094683}{Sequence record}, accessed 9 September
  2026.
\item
  S. W. Graham and G. Kolesnik, \emph{Van der Corput's Method of
  Exponential Sums}, London Mathematical Society Lecture Note Series
  126, Cambridge University Press, 1991.
  \href{https://beckassets.blob.core.windows.net/product/readingsample/666252/9780521339278_excerpt_001.pdf}{Publisher
  excerpt}.
\item
  L. Kuipers and H. Niederreiter, \emph{Uniform Distribution of
  Sequences}, Wiley-Interscience, New York, 1974. Chapter 2, especially
  pp.~112--116.
  \href{https://web.maths.unsw.edu.au/~josefdick/preprints/KuipersNied_book.pdf}{University-hosted
  scan}.
\item
  R. Terras, ``A stopping time problem on the positive integers,''
  \emph{Acta Arithmetica} 30 (1976), 241--252.
  \href{https://doi.org/10.4064/aa-30-3-241-252}{doi:10.4064/aa-30-3-241-252}.
\end{enumerate}
\end{document}
